% Resumo em língua inglesa
\setlength{\absparsep}{18pt}
\begin{resumo}[Abstract]
 \begin{otherlanguage*}{english}
Adverse Events (AEs) in hospital environments represent one of the greatest contemporary global public health challenges, ranking among the leading preventable causes of inpatient morbidity and mortality worldwide. Historically, hospital surveillance systems have relied almost exclusively on voluntary and spontaneous incident reports, an inherently flawed mechanism that captures less than 10\% of events with actual patient harm, thereby reinforcing a false sense of institutional safety. Conversely, the Global Trigger Tool (GTT) methodology, established by the Institute for Healthcare Improvement (IHI), has emerged internationally as the gold standard for retrospective AE audit through the identification of clinical triggers and standardized harm severity categorization defined by the National Coordinating Council for Medication Error Reporting and Prevention (NCC MERP). Nevertheless, formal education in patient safety within undergraduate healthcare curricula faces substantial bioethical barriers, legal liabilities, and data privacy restrictions imposed by regulations such as Brazil's General Data Protection Law (LGPD), precluding unrestricted student access to authentic patient records. This departmental project formalizes the software engineering specification and academic foundation of \textbf{SIGEA-GTT} (Intelligent Adverse Event Management System – Global Trigger Tool), an innovative web-based platform tailored for training, simulation, and clinical auditing at the Federal University of Sergipe (UFS). The software ensures rigorous sensitive data isolation by providing structured simulated electronic health records (multidisciplinary progress notes, scheduled medication orders, laboratory tests, and surgical procedures), a standardized 20-minute audit timer, and screening across 53 canonical IHI triggers divided into 6 care modules. As a core pedagogical differentiator, SIGEA-GTT natively incorporates six Quality Improvement Tools (Ishikawa 6M Diagram, real-time GUT Matrix, 5W2H Action Plan, PDCA/PDSA Cycle, SWOT Matrix, and Brainstorming), formative instructor feedback, and automated computation of the three official IHI epidemiological metrics: AEs per 1,000 patient-days, AEs per 100 admissions, and percentage of admissions with adverse events. The platform is architected upon Java 21 LTS with Spring Boot 3.3, Angular 18+ featuring Standalone components and Signals reactivity, relational PostgreSQL 16 persistence utilizing JSONB attributes, and stateless JWT authentication restricted to the institutional academic domain. The findings establish a pioneering educational framework in Brazil for rigorous patient safety instruction driven by high-reliability software engineering principles.

   \vspace{\onelineskip}
 
   \noindent 
   \textbf{Keywords}: Patient Safety. Global Trigger Tool. Adverse Events. Quality Improvement Tools. Computational Clinical Simulation. Software Engineering.
 \end{otherlanguage*}
\end{resumo}